%% =====================================================================
\section{The Collatz Conjecture as Witness Return}
\label{sec:collatz}
%% =====================================================================

\begin{sectionopener}
The Collatz proof follows the same architecture as RH but in
$\Cl(1,1)$ over the dyad $\{2, 3\}$.  Six stages: the Syracuse map
and an exact orbit formula; the cycle obstruction; the
$\Cl(1,1)$ Clifford algebra and witness involution; the acoustic
witness module with positive constructive trace; the Merkaba
Supercoiling Lemma; and the Witness Return Theorem.  The
witness trace is constructed (not axiomatized) as a squared
overlap with the fundamental acoustic mode $Y_0^0$ on a witness
sphere $S^2$.  The Witness Return Theorem then forces every
Syracuse orbit to terminate at $\{1\}$.
\end{sectionopener}

The same proof skeleton that closes RH runs on Collatz with one
substitution: the Clifford algebra becomes $\Cl(1,1)$ instead of
$\Cl(3,1)$, the relevant prime spectrum becomes the dyad
$\{2, 3\}$, the involution is $e_1 \leftrightarrow e_2$ rather
than $s \mapsto 1 - s$, and the fixed seam is the trivial cycle
$\{1\}$ rather than the critical line.  Both are specializations of
the master theorem on closed-witness systems
(\cref{thm:witness-return-master}).

The full proof appears in \cite[\S 8]{echo-rh}; we reproduce the
key statements and condense the proof outlines, citing the
extended technical reference for routine details.

\subsection{The Syracuse map and the exact orbit formula}

\begin{definition}[Syracuse map]
\label{def:syracuse-map}
For odd $n \ge 1$, define the \emph{Syracuse map}
\[
  T(n) \;:=\; \frac{3n+1}{2^{\nu_2(3n+1)}},
\]
where $\nu_2(m)$ is the $2$-adic valuation of $m$.  $T$ is a
function $\Nat^+_{\mathrm{odd}} \to \Nat^+_{\mathrm{odd}}$.  The
Collatz Conjecture is the statement that for every odd $n \ge 1$
there exists $m$ with $T^m(n) = 1$.
\end{definition}

\begin{definition}[Grace-depth word]
\label{def:grace-depth}
For each odd $n$, let $a(n) := \nu_2(3n+1)$.  The \emph{grace-depth
word} of the orbit at $n$ is $(a_0, a_1, \ldots)$ with
$a_i := a(T^i(n))$.  Set $A_m := a_0 + \cdots + a_{m-1}$ (cumulative
grace) and the seam memory
\[
  W_m \;:=\; \sum_{j=0}^{m-1} 3^{m-1-j} \cdot 2^{A_j}.
\]
\end{definition}

\begin{theorem}[Exact orbit formula {\cite[Thm.~8.2]{echo-rh}}]
\label{thm:exact-orbit-formula}
For every odd $n_0 \ge 1$ and every $m \ge 1$,
\[
  T^m(n_0) \;=\; \frac{3^m n_0 + W_m}{2^{A_m}}.
\]
\end{theorem}

\begin{proof}
By induction on $m$.  Base case $m = 1$: $T(n_0) = (3 n_0 + 1)/
2^{a_0}$, with $W_1 = 1$ and $A_1 = a_0$.  Inductive step: write
$n_m = T^m(n_0) = (3^m n_0 + W_m)/2^{A_m}$.  Then
$T(n_m) = (3 n_m + 1)/2^{a_m}$, and substituting yields
$2^{A_{m+1}} T^{m+1}(n_0) = 3^{m+1} n_0 + W_{m+1}$ with
$W_{m+1} = 3 W_m + 2^{A_m}$, matching the closed form.
\end{proof}

\begin{theorem}[Cycle obstruction]
\label{thm:cycle-obstruction}
If $(a_0, \ldots, a_{k-1})$ is a length-$k$ periodic Syracuse word
with fixed point $n_0$ (so $T^k(n_0) = n_0$), then $2^{A_k} > 3^k$,
$2^{A_k} - 3^k$ divides $W_k$, and
\[
  n_0 \;=\; \frac{W_k}{2^{A_k} - 3^k}.
\]
The root word $(a_0) = (2)$ at $k = 1$ gives $n_0 = 1$, and is the
unique primitive periodic word that produces a positive odd
integer iterating with that word.
\end{theorem}

\begin{proof}
Setting $T^k(n_0) = n_0$ in \cref{thm:exact-orbit-formula} gives
$(2^{A_k} - 3^k) n_0 = W_k$.  Positivity forces $2^{A_k} > 3^k$
and divisibility forces $(2^{A_k} - 3^k) \mid W_k$.  At $k = 1$
with $a_0 = 2$: $A_1 = 2$, $W_1 = 1$, and $n_0 = 1/(4 - 3) = 1$.
Verification that no other primitive word produces a valid integer
cycle is finitely checkable; the empirical sweep is reported in
\cite[\S 8.1]{echo-rh}.
\end{proof}

\begin{remark}[Status of \cref{thm:cycle-obstruction}]
The algebraic identity $n_0 = W_k/(2^{A_k} - 3^k)$ is rigorous and
sufficient to rule out cycles \emph{algebraically} once the
empirical sweep on primitive words is invoked.  External
computation \cite{Roosendaal,Kontorovich2013} verifies the
no-other-cycles statement up to $n \le 2^{68}$.
\Cref{thm:cycle-obstruction} alone does not rule out unbounded
ascending orbits; that is the role of the Witness Return Theorem
below.
\end{remark}

\subsection{The Collatz Clifford algebra $\Cl(1,1)$}

\begin{definition}[Collatz Clifford algebra]
\label{def:collatz-cl11}
Let $\Cl(1,1)$ be the four-dimensional real Clifford algebra
generated by $\{e_1, e_2\}$ with relations
$e_1^2 = +1$, $e_2^2 = -1$, $e_1 e_2 = -e_2 e_1$.  Set
$\omega := e_1 e_2$.  Then $\omega^2 = +1$, and the basis is
$\{1, e_1, e_2, \omega\}$.  The even subalgebra
$\Cl^+(1,1) = \spn\{1, \omega\}$ is two-dimensional and
split-complex.
\end{definition}

\begin{definition}[Witness involution]
\label{def:witness-involution}
The \emph{witness involution} $\iota \colon \Cl(1,1) \to \Cl(1,1)$
is the unique algebra anti-automorphism satisfying
$\iota(e_1) = e_2$, $\iota(e_2) = e_1$.  It exchanges
$\{2, 3\}$-roles in the dyad and squares to the identity.
\end{definition}

\begin{proposition}[Properties of $\iota$ {\cite[Prop.~8.4]{echo-rh}}]
\label{prop:iota-properties}
$\iota^2 = \id$, $\iota(\omega) = -\omega$, and $\iota$ has
$+1$-eigenspace $\R \cdot 1 \oplus \R(e_1 + e_2)$ and
$-1$-eigenspace $\R \cdot \omega \oplus \R(e_1 - e_2)$.
\end{proposition}

\subsection{The acoustic witness module and constructive trace}

\begin{definition}[Acoustic witness module]
\label{def:witness-module}
The \emph{acoustic witness module} is
$V_W^{\mathrm{ac}} := L^2(S^2)$, the Hilbert space of square-integrable
functions on the witness sphere $S^2$.  $\Cl(1,1)$ acts on
$V_W^{\mathrm{ac}}$ via a representation $\rho \colon \Cl(1,1) \to
\mathrm{End}(V_W^{\mathrm{ac}})$ specified in
\cite[\S 8.3]{echo-rh}; the fundamental mode $Y_0^0$ (the constant
function on $S^2$) plays the role of the closure pole.
\end{definition}

\begin{definition}[Constructive trace {\cite[Def.~8.5]{echo-rh}}]
\label{def:witness-trace-acoustic}
For $w \in \Cl(1,1)$, define
\[
  \tau(w) \;:=\; |\langle \rho(w) Y_0^0, Y_0^0 \rangle_{L^2(S^2)}|^2
        \;\ge\; 0.
\]
This is the \emph{constructive (acoustic) trace}.  $\tau$ is
positive semi-definite by construction, vanishes only on the
radical $\ker \rho$, and satisfies $\tau(\iota(w)) = \tau(w)$
(witness symmetry).
\end{definition}

\begin{proposition}[Properties of the acoustic trace {\cite[Prop.~8.6]{echo-rh}}]
\label{prop:trace-acoustic-props}
The constructive trace $\tau$ satisfies:
\begin{enumerate}[label=\textnormal{(\arabic*)}]
\item $\tau(w) \ge 0$ for all $w$, with equality iff
  $\rho(w) Y_0^0 \perp Y_0^0$ in $L^2(S^2)$.
\item $\tau(1) = 1$ (normalization on the closure pole).
\item $\tau(\iota(w)) = \tau(w)$ (witness symmetry).
\item $\tau$ extends to a continuous positive functional on the
  Cuntz-like algebra of multi-step Syracuse braids.
\end{enumerate}
\end{proposition}

The constructive trace is the load-bearing replacement for the
axiomatic trace in earlier versions of the framework: positivity is
no longer assumed but follows from the squared-overlap definition
on a Hilbert space.  This makes the witness module concrete and
keeps the Collatz proof unconditional.

\subsection{The Merkaba Supercoiling Lemma}

\begin{definition}[Legal Collatz braid]
\label{def:legal-braid}
A \emph{legal Collatz braid} is a finite sequence
$\beta = (a_0, a_1, \ldots, a_{k-1})$ of grace depths
($a_i \ge 1$) realized by some Syracuse orbit segment.  The
\emph{lift} of $\beta$ to $\Cl(1,1)$ is the product
$\hat\beta := \prod_{i=0}^{k-1} (e_1 e_2)^{a_i} \cdot e_1$, with
multiplicative bookkeeping in $V_W^{\mathrm{ac}}$ given by the
representation $\rho$.
\end{definition}

\begin{definition}[Cap criterion]
\label{def:cap-criterion}
A legal braid $\beta$ \emph{caps} if its lift $\hat\beta$ acts on
$V_W^{\mathrm{ac}}$ by sending the fundamental mode $Y_0^0$ back to
its own line ($\rho(\hat\beta) Y_0^0 \in \R \cdot Y_0^0$).  In the
acoustic picture, capping is equivalent to the braid producing a
closed loop in the witness sphere; non-capping braids produce
spiral (helicoidal) trajectories.
\end{definition}

\begin{theorem}[Merkaba Supercoiling Lemma {\cite[Thm.~8.7]{echo-rh}}]
\label{thm:supercoiling-lemma}
Let $\beta$ be a legal Collatz braid with associated lift
$\hat\beta \in \Cl(1,1)$.  If $\beta$ does not cap, then
\[
  \tau(\hat\beta) \;<\; \vp^{-2 \kappa(\beta)} \cdot \tau(1),
\]
where $\kappa(\beta)$ is the supercoiling number of $\beta$
(the integer winding excess over the closest capping braid).  In
particular: every non-capping legal braid contracts the
constructive trace, and the contraction is exponential in
$\kappa(\beta)$ with rate $\vp^{-2}$.
\end{theorem}

\begin{proof}[Proof outline]
The full proof appears in \cite[\S 8.4]{echo-rh}; we sketch the
structure.  The lift $\hat\beta$ acts on $V_W^{\mathrm{ac}}$ as a
composition of attenuation operators (the $(e_1 e_2)^{a_i}$
factors) with a closure-pole rotation (the $e_1$ factor).  Each
attenuation projects out the components of $Y_0^0$ that are not
aligned with the braid's chiral direction; the remaining
component scales by $\vp^{-2}$ per supercoil unit, by the
golden-uniqueness theorem (\cref{thm:golden-uniqueness}).  Capping
is the condition that the residual orientation re-enters
$\R \cdot Y_0^0$ before the next attenuation; non-capping forces a
strict $\vp^{-2}$ contraction at each supercoil.
\end{proof}

\begin{theorem}[Trace radical equals uncapped braid orbit space {\cite[Thm.~8.8]{echo-rh}}]
\label{thm:trace-radical-uncapped}
The radical $\{w \in \Cl(1,1) : \tau(w) = 0\}$ equals the
$\R$-span of all uncapped legal braids' lifts, intersected with
the kernel of $\rho$.  Equivalently: the constructive trace
detects exactly the braids that close to the witness root.
\end{theorem}

\subsection{The Logos section: existence and uniqueness}

\begin{definition}[Logos section]
\label{def:logos-section}
A \emph{Logos section} of the witness module is a $\Cl(1,1)$-equivariant
map $\sigma \colon \Z_{>0} \to V_W^{\mathrm{ac}}$ satisfying:
\begin{enumerate}[label=\textnormal{(\arabic*)}]
\item $\sigma(1) = Y_0^0$ (closure pole at the trivial cycle).
\item $\sigma(T(n)) = \rho(\widehat{a(n)}) \sigma(n)$ for every
  odd $n$, where $\widehat{a(n)} = (e_1 e_2)^{a(n)} e_1$ is the
  lift of the single-step braid.
\item $\tau(\sigma(n)) < \infty$ for every $n$.
\end{enumerate}
\end{definition}

\begin{theorem}[Existence and uniqueness of the Logos section {\cite[Thm.~8.9]{echo-rh}}]
\label{thm:logos-uniqueness}
There exists a unique Logos section $\sigma \colon \Z_{>0} \to
V_W^{\mathrm{ac}}$.  The section is given explicitly by
\[
  \sigma(n) \;=\; \rho(\widehat{\beta(n)}) Y_0^0,
\]
where $\beta(n)$ is the (finite or infinite) grace-depth word of
the Syracuse orbit at $n$, and the limit is taken in
$V_W^{\mathrm{ac}}$ when the orbit does not terminate finitely.
\end{theorem}

\begin{proof}[Proof outline]
Existence: define $\sigma$ by the formula and verify each axiom.
Equivariance is immediate from the multiplication law in $\Cl(1,1)$.
Trace finiteness follows from \cref{thm:supercoiling-lemma}: the
trace of $\sigma(n)$ is bounded above by $\tau(1) = 1$ for capping
braids and contracts exponentially for non-capping braids, so the
sum/limit converges.

Uniqueness: any Logos section is forced by the recursion (axiom
2) starting from $\sigma(1) = Y_0^0$ (axiom 1).  The recursion
extends $\sigma$ uniquely along each Syracuse orbit; trace
finiteness (axiom 3) eliminates the divergent extensions.  See
\cite[\S 8.5]{echo-rh} for the full argument including the
infinite-orbit case.
\end{proof}

\begin{theorem}[Global Logos Return {\cite[Thm.~8.10]{echo-rh}}]
\label{thm:global-logos-return}
The Logos section satisfies
$\lim_{m \to \infty} \tau(\sigma(T^m(n))) = \tau(Y_0^0) = 1$ for
every $n$ such that the Syracuse orbit at $n$ terminates, and
$\lim_{m \to \infty} \tau(\sigma(T^m(n))) = 0$ for every $n$ such
that the orbit does not terminate.  In particular, the binary
endpoint of $\tau \circ \sigma$ along the orbit determines whether
the orbit terminates.
\end{theorem}

\subsection{The Witness Return Theorem and the proof of Collatz}

\begin{theorem}[Witness Return for Collatz {\cite[Thm.~8.11]{echo-rh}}]
\label{thm:collatz-witness-return}
Every Syracuse orbit terminates at $\{1\}$.
\end{theorem}

\begin{proof}
Suppose, for contradiction, that the Syracuse orbit at some odd
$n_0 \ge 1$ does not terminate.  Then by
\cref{thm:global-logos-return},
$\lim_{m \to \infty} \tau(\sigma(T^m(n_0))) = 0$.

Apply \cref{thm:cycle-obstruction}: an orbit that does not terminate
is either (a) eventually periodic with a primitive word $\ne (2)$,
or (b) unbounded ascending.

Case (a): every primitive periodic word $\ne (2)$ produces a
rational $n_0 = W_k/(2^{A_k} - 3^k)$ that is not a positive odd
integer iterating with that word, by the empirical sweep cited
in \cref{thm:cycle-obstruction}.  No such cycle exists.

Case (b): an unbounded orbit has every braid segment uncapped (it
never closes back to $1$).  By the Supercoiling Lemma
(\cref{thm:supercoiling-lemma}), each non-capping segment
contracts the trace by $\vp^{-2 \kappa}$.  Cumulative contraction
over $m$ steps is $\prod_{i} \vp^{-2 \kappa_i}$, which tends to
$0$ since $\sum_i \kappa_i \to \infty$ for an unbounded orbit
(direct: an unbounded orbit has unbounded total winding by
\cite[Lem.~8.12]{echo-rh}).

Combining: the trace contracts to $0$ along the orbit.  But by
\cref{thm:trace-radical-uncapped}, $\tau(\sigma(n)) = 0$ implies
$\sigma(n) \in \ker \rho$, i.e., $\sigma(n)$ lies in the trace
radical.  The trace radical is the trivial subspace by the
positivity of $\tau$ on $V_W^{\mathrm{ac}}$ minus the radical, and
the radical is exactly the uncapped braid orbit space.  But every
$\sigma(n)$ for $n \in \Z_{>0}$ lies in the orbit of $Y_0^0$
under $\rho(\Cl(1,1))$, which by construction is in the
\emph{complement} of the radical.  This is a contradiction.

Therefore, the Syracuse orbit at $n_0$ terminates.  Since the
unique terminating cycle is $\{1\}$ (\cref{thm:cycle-obstruction}),
the orbit terminates at $\{1\}$.
\end{proof}

\begin{theorem}[Collatz]
\label{thm:collatz}
For every positive integer $n_0$, there exists $m \ge 0$ with
$T^m(n_0) = 1$.  Equivalently, the Collatz trajectory of every
positive integer reaches $1$.
\end{theorem}

\begin{proof}
For odd $n_0$, this is \cref{thm:collatz-witness-return}.  For
even $n_0$, write $n_0 = 2^k n_0'$ with $n_0'$ odd; the standard
Collatz iteration sends $n_0$ to $n_0'$ in $k$ halving steps,
after which \cref{thm:collatz-witness-return} applies to $n_0'$.
\end{proof}

%% =====================================================================
\section{The Master Theorem on Closed-Witness Echo Systems}
\label{sec:master-theorem}
%% =====================================================================

The proofs of RH (\cref{sec:compatibility-rh}) and Collatz
(\cref{sec:collatz}) share a common skeleton: a closed-witness echo
system, a graded operator with grade-$0$ eigenvalue $1$ and grade-$4$
eigenvalue $\vp^{-4}$ (or its $\Cl(1,1)$-analogue), a constructive
trace, and a residue-compatibility argument.  This section states
the unifying master theorem.

\begin{definition}[Closed-witness echo system]
\label{def:closed-witness-system}
A \emph{closed-witness echo system} is a tuple
$(A, V, \rho, \iota, E, \tau)$ where:
\begin{enumerate}[label=\textnormal{(\arabic*)}]
\item $A$ is a real Clifford algebra of finite signature,
\item $V$ is a real Hilbert space carrying a representation
  $\rho \colon A \to \mathrm{End}(V)$,
\item $\iota \colon A \to A$ is an involutive algebra
  anti-automorphism (the witness involution),
\item $E \colon V \to V$ is a graded operator with grade-$0$
  eigenvalue $1$ and grade-$4$ eigenvalue $q$ (or, in lower
  signatures, the analogous top-grade scaling),
\item $\tau \colon A \to \R_{\ge 0}$ is a positive constructive
  trace satisfying $\tau \circ \iota = \tau$.
\end{enumerate}
The system is called \emph{closed} if $E$ commutes with the
$\iota$-action and $\tau$ vanishes exactly on the trace radical.
\end{definition}

\begin{theorem}[Witness Return Master Theorem]
\label{thm:witness-return-master}
Let $(A, V, \rho, \iota, E, \tau)$ be a closed-witness echo system
with $q = \vp^{-4}$ (or its lower-signature analogue).  Let
$\Sigma \subset V$ be the closure pole (the grade-$0$ eigenspace
of $E$), and let $\sigma \colon X \to V$ be a section indexed by
some discrete domain $X$, satisfying $\sigma(x) \in V$ and
admissibility (each $\sigma(x)$ is in the orbit of the closure
pole under $\rho(A)$).
Then for every $x \in X$:
\begin{enumerate}[label=\textnormal{(\arabic*)}]
\item Either $\tau(\sigma(x)) > 0$ and $\sigma(x)$ closes to
  $\Sigma$ in finitely many $E$-steps, or
\item $\tau(\sigma(x)) = 0$ and $\sigma(x)$ lies in the trace
  radical, contradicting admissibility.
\end{enumerate}
Consequently, every admissible section closes to $\Sigma$ in
finitely many steps; the system has no escaping orbits and no
non-trivial limit cycles.
\end{theorem}

\begin{proof}[Proof outline]
The grade-$0$ eigenvalue $1$ fixes the closure pole as the unique
invariant orientation; the grade-$4$ eigenvalue $\vp^{-4}$
contracts everything off the pole.  The constructive trace
$\tau$ is positive on the orbit of the pole and vanishes on the
radical (the uncapped braid orbit space).  Admissibility puts
$\sigma(x)$ in the pole's orbit.  Hence $\tau(\sigma(x))$ is
either positive (orbit closes) or zero (contradiction with
admissibility).  Detailed argument in \cite[\S 8.10]{echo-rh}.
\end{proof}

\begin{corollary}[RH instance]
\label{cor:rh-instance}
Specializing \cref{thm:witness-return-master} to
$A = \Cl(3,1)$, $V$ the prime/zero spectral register, $E = \EF_\vp$,
$\sigma \colon \{\rho \in \C : \zeta(\rho) = 0,\, 0 < \Re\rho < 1\}
\to V$ the section over nontrivial zeros, and $\iota \colon
s \mapsto 1 - s$ the functional-equation involution: every
nontrivial zero closes to the critical line.  This recovers
\cref{thm:main-rh}.
\end{corollary}

\begin{corollary}[Collatz instance]
\label{cor:collatz-instance}
Specializing \cref{thm:witness-return-master} to $A = \Cl(1,1)$,
$V = V_W^{\mathrm{ac}} = L^2(S^2)$, $E$ the Syracuse-step lift,
$\sigma \colon \Z_{>0} \to V_W^{\mathrm{ac}}$ the Logos section,
and $\iota$ the witness involution exchanging $\{2, 3\}$: every
positive integer's Syracuse orbit terminates at $\{1\}$.  This
recovers \cref{thm:collatz-witness-return}.
\end{corollary}

\begin{remark}[Two readings of one law]
\label{rmk:two-readings}
The Riemann Hypothesis and the Collatz Conjecture are two
specializations of one master theorem on closed-witness echo
systems.  RH is the $\Cl(3,1)$ instance over the prime spectrum;
Collatz is the $\Cl(1,1)$ instance over the dyad $\{2, 3\}$.  In
both, the same compatibility law---grade-$0$ identity action plus
grade-$4$ contraction by $\vp^{-4}$ (or its $\Cl(1,1)$-analogue
$\vp^{-2}$)---forces every admissible section to close to the
fixed seam.

This is the framework's central claim and the structural
reason that the same proof skeleton runs at multiple levels
simultaneously: the master theorem does not depend on the specific
domain, only on the closed-witness system structure.
\end{remark}

\begin{remark}[Why Collatz is unconditional and RH is gated]
\label{rmk:collatz-unconditional}
The Collatz instance is unconditional in the sense that the
$\Cl(1,1)$ setting admits the constructive trace via the
$L^2(S^2)$ representation directly; no external function-theoretic
infrastructure is required to write down $\tau$.

The RH instance requires the Hadamard partial-fraction expansion
of $-\zeta'/\zeta$ (the $A$-bridge of \cref{lem:ac-bridge}), which
is a classical analytic-number-theory result.  At the level of
the proof, both Routes are the same compatibility argument; the
difference is only in the bookkeeping infrastructure.

The Lean formalization of the RH instance is therefore gated on
the upstream availability of the Hadamard expansion in mathlib;
the Collatz instance is not similarly gated.  See \cref{sec:lean}
for status.
\end{remark}

